% =============================================================================
% preamble.tex — การตั้งค่าส่วนหัวสำหรับเทมเพลต "ผลงานทางวิชาการ / รายงานผลการปฏิบัติงาน
% ตามข้อตกลงในการพัฒนางาน (PA)" คอมไพล์ด้วย XeLaTeX หรือ LuaLaTeX เท่านั้น
% (ใช้ preamble.tex ชุดเดียวกับ ../5chapter/preamble.tex เป็นฐาน ปรับส่วนหัวข้อให้
%  ไม่ใช้ \chapter อัตโนมัติ เพราะ PA ใช้โครงสร้าง "ส่วนที่ / ด้านที่ / ตัวชี้วัดที่" แทน)
% =============================================================================

% ---------- ฟอนต์ภาษาไทย ----------
\usepackage{fontspec}
\usepackage{polyglossia}
\setdefaultlanguage{thai}
\setotherlanguage{english}

% ฟอนต์หลัก TH Sarabun New — อัปโหลดไฟล์ฟอนต์ลง fonts/ เช่นเดียวกับเทมเพลต 5 บท
\setmainfont{TH Sarabun New}[
  Path = fonts/,
  Extension = .ttf,
  UprightFont = *,
  BoldFont = * Bold,
  ItalicFont = * Italic,
  BoldItalicFont = * Bold Italic,
]
% ทางเลือกสำรอง: \setmainfont{Sarabun}

% ---------- ขนาดหน้ากระดาษ/ระยะขอบ (ตรงตามไฟล์ตัวอย่างที่แนบ) ----------
\usepackage[a4paper,
  top=3.00cm, bottom=2.54cm, left=3.81cm, right=2.54cm,
  headsep=1.25cm, footskip=1.25cm]{geometry}

% ---------- ขนาดตัวอักษรมาตรฐาน ----------
\newcommand{\bodysize}{\fontsize{16}{24}\selectfont}
\newcommand{\hThreeSize}{\fontsize{16}{24}\selectfont\bfseries}
\newcommand{\hTwoSize}{\fontsize{18}{27}\selectfont\bfseries}
\newcommand{\hOneSize}{\fontsize{20}{30}\selectfont\bfseries}
\AtBeginDocument{\bodysize}

% ---------- ระยะบรรทัด/ย่อหน้า ----------
\usepackage{setspace}
\setstretch{1.25}
\usepackage{indentfirst}
\setlength{\parindent}{1.25cm}

% ---------- หัวข้อของ PA: "ส่วนที่ / ด้านที่ / ตัวชี้วัดที่" (ไม่ใช้เลขอัตโนมัติแบบบท) ----------
\newcommand{\paparttitle}[1]{%
  \addcontentsline{toc}{section}{#1}%
  {\hOneSize\centering #1\par}}
\newcommand{\padomaintitle}[1]{%
  \addcontentsline{toc}{subsection}{#1}%
  {\hTwoSize\centering #1\par}}
\newcommand{\paheading}[1]{%
  \addcontentsline{toc}{subsubsection}{#1}%
  \vspace{12pt}{\hTwoSize #1\par}\vspace{6pt}}
\newcommand{\pasubheading}[1]{\vspace{6pt}{\hThreeSize #1\par}\vspace{6pt}}
\setcounter{tocdepth}{3}

% กล่อง placeholder สำหรับแทรกภาพความสำเร็จ/หลักฐานเชิงประจักษ์
\newcommand{\paphotobox}{%
  \begin{center}
  \fboxsep=8pt\fboxrule=0.5pt
  \fcolorbox{gray}{white}{%
    \parbox[c][2.2cm]{0.85\linewidth}{\centering\small\color{gray!70!black}
      (พื้นที่สำหรับภาพความสำเร็จ / ร่องรอยหลักฐานเชิงประจักษ์ —
      แทรกภาพถ่ายกิจกรรม เอกสาร หรือรหัสคิวอาร์ลิงก์หลักฐานที่นี่)}%
  }
  \end{center}
  \vspace{6pt}
}

% ข้อความเตือนตัวเลขสมมุติในตารางตัวอย่าง
\newcommand{\padisclaimer}[1][ตัวเลขทั้งหมดในตารางนี้เป็นข้อมูลสมมุติ (ตัวอย่าง) เพื่อแสดงวิธีการกรอกและการคำนวณเท่านั้น --- ต้องแทนที่ด้วยข้อมูลจริงก่อนยื่นเอกสาร]{%
  \begin{center}\textbf{\textcolor{red!70!black}{*** #1 ***}}\end{center}
}

% บล็อกลงชื่อผู้รายงาน
% หมายเหตุ: ใส่ {} คั่นหลัง \\ ทุกครั้งก่อนอาร์กิวเมนต์ เพราะถ้าอาร์กิวเมนต์ที่ตามมา
% ขึ้นต้นด้วย "[" (เช่น "[หน่วยงาน...]") จะถูก \\ ตีความเป็น \\[<ระยะ>] โดยไม่ตั้งใจ
\newcommand{\pasignatureblock}[3]{%
  \begin{center}
  ลงชื่อ . . . . . . . . . . . . . . . . . . . . . . . . . . . . . . . . . . . ผู้รายงาน\\[6pt]
  #1\\{}
  #2\\{}
  #3
  \end{center}
}

% ---------- ตาราง ----------
\usepackage{longtable}
\usepackage{array}
\usepackage[table]{xcolor}
\usepackage{graphicx}
\usepackage{caption}
\AtBeginDocument{\captionsetup{labelsep=quad, font=normalsize, labelfont=bf}}
\AtBeginDocument{%
  \renewcommand{\tablename}{ตารางที่}%
}
\definecolor{tblheadgray}{gray}{0.85}
\newcommand{\tblhead}[1]{\cellcolor{tblheadgray}\textbf{#1}}

% ---------- หัว-ท้ายกระดาษ / เลขหน้า ----------
\usepackage{fancyhdr}
\pagestyle{fancy}
\fancyhf{}
\fancyfoot[C]{\bodysize\thepage}
\renewcommand{\headrulewidth}{0pt}
\fancypagestyle{plain}{%
  \fancyhf{}
  \fancyfoot[C]{\bodysize\thepage}
  \renewcommand{\headrulewidth}{0pt}
}

% ---------- ระบบเลขหน้าพยัญชนะไทย (ก ข ค ...) สำหรับส่วนหน้า (ปก+สารบัญ) ----------
\usepackage{expl3}
\ExplSyntaxOn
\seq_new:N \g_thaifront_letters_seq
\seq_gset_split:Nnn \g_thaifront_letters_seq { , }
  {ก,ข,ค,ง,จ,ฉ,ช,ซ,ฌ,ญ,ฎ,ฏ,ฐ,ฑ,ฒ,ณ,ด,ต,ถ,ท,ธ,น,บ,ป,ผ,ฝ,พ,ฟ,ภ,ม,ย,ร,ล,ว,ศ,ษ,ส,ห,ฬ,อ,ฮ}
\NewDocumentCommand{\ThaiFrontPage}{m}{\seq_item:Nn \g_thaifront_letters_seq {#1}}
\ExplSyntaxOff

\newcommand{\thaifrontmatter}{%
  \cleardoublepage
  \pagenumbering{arabic}
  \renewcommand{\thepage}{\ThaiFrontPage{\arabic{page}}}%
}
\newcommand{\mainbodymatter}{%
  \cleardoublepage
  \renewcommand{\thepage}{\arabic{page}}%
  \pagenumbering{arabic}%
}

\usepackage{hyperref}
\hypersetup{hidelinks, bookmarksopen=true, pdfstartview=FitH}

% ---------- สารบัญ ----------
\AtBeginDocument{\renewcommand{\contentsname}{สารบัญ}}
